Given the \texttt{Employee} class \textbf{provided in the template}, extend it to create three new subclasses by overriding the necessary methods:

\begin{itemize}

    \item \texttt{Lawyer}:

        \begin{itemize}

            \item Has half the vacation days of a regular \texttt{Employee}.

            \item Earns 10000 euros more than an \texttt{Employee}.

            \item Works 5 more hours per week than an \texttt{Employee}.

            \item Takes vacation in July instead of August.

            \item Has an additional method called \texttt{attendCourt()} that prints: \texttt{"I am in court"}.

        \end{itemize}

    \item \texttt{Secretary}:

        \begin{itemize}

            \item Earns 5000 euros more than an \texttt{Employee}.

            \item Works 30 hours per week, same as an \texttt{Employee}.

            \item Has the same vacation days as an \texttt{Employee} and takes them in the same month.

            \item Has an additional method called \texttt{makePhotocopies()} that prints: \texttt{"I am making photocopies"}.

        \end{itemize}

    \item \texttt{Janitor}:

        \begin{itemize}

            \item Has 5 more vacation days than an \texttt{Employee}.

            \item Takes vacation in September instead of August.

            \item Works the same hours and earns the same salary as an \texttt{Employee}.

        \end{itemize}

\end{itemize}

Create a \texttt{Main} class that:

\begin{itemize}

    \item Stores ten employees in an array, mixing the four types in any order.

    \item Iterates through the array, calling and displaying the output of all the defined methods for each employee. \textit{(Hint: You will need to check the specific type of the object to call the additional methods).}

\end{itemize}
